\documentclass{article}

\usepackage[T1]{fontenc}
\usepackage[utf8]{inputenc}
\usepackage{multicol}
\usepackage{calc}
\usepackage{ifthen}
\usepackage{epsfig}
%\usepackage{bookman}
%\usepackage{booktabs}
\usepackage{upquote} % Uncurly the quotes
\usepackage{etoolbox} % for backquote fix
\usepackage[scaled=0.8]{DejaVuSans}
\usepackage[scaled=0.8]{DejaVuSansMono}
\usepackage{mdframed} % nice frames
\usepackage[nobottomtitles]{titlesec} % better titles
\usepackage{enumitem}
\usepackage{minted}
\usepackage[a4paper,landscape,margin=1cm,includehead,headsep=2ex,nofoot]{geometry}
\usepackage{fancyhdr}\usepackage{hyperref}

\usemintedstyle{tango}

\setlist[itemize]{noitemsep, topsep=0pt}

% backquote dejavu fix
\makeatletter
\chardef\straightquote@code=\catcode`'
\chardef\backquote@code=\catcode``
\catcode`'=\active \catcode``=\active
\patchcmd{\@noligs}
 {\textasciigrave}
 {\fixedtextasciigrave}
 {}{}
\newcommand{\fixedtextasciigrave}{%
  \makebox[.5em]{\fontencoding{TS1}\fontfamily{fvs}\selectfont\textasciigrave}% Vera Sans
}
\catcode\lq\'=\straightquote@code
\catcode\lq\`=\backquote@code
\makeatletter

% fancyvrb and mdframed stuff
\newmdenv[bottomline=false,rightline=false]{topleft}
\newmdenv[topline=false,leftline=false]{bottomright}
\DefineVerbatimEnvironment{Verbacorner}{Verbatim}{}
\surroundwithmdframed[bottomline=false,rightline=false]{Verbacorner}
\fvset{commandchars=@\{\}}

\ifthenelse{\lengthtest { \paperwidth = 11in}}
	{ \geometry{top=.5in,left=.5in,right=.5in,bottom=.5in} }
	{\ifthenelse{ \lengthtest{ \paperwidth = 297mm}}
		{\geometry{top=1cm,left=1cm,right=1cm,bottom=1cm} }
		{\geometry{top=1cm,left=1cm,right=1cm,bottom=1cm} }
	}
\pagestyle{empty}

%% \makeatletter
%% \renewcommand{\section}{\@startsection{section}{1}{0mm}%
%%                                 {-1ex plus -.5ex minus -.2ex}%
%%                                 {0.5ex plus .2ex}%x
%%                                 {\normalfont\large\bfseries}}
%% \renewcommand{\subsection}{\@startsection{subsection}{2}{0mm}%
%%                                 {-1ex plus -.5ex minus -.2ex}%
%%                                 {0.5ex plus .2ex}%
%%                                 {\normalfont\normalsize\bfseries}}
%% \renewcommand{\subsubsection}{\@startsection{subsubsection}{3}{0mm}%
%%                                 {-1ex plus -.5ex minus -.2ex}%
%%                                 {1ex plus .2ex}%
%%                                 {\normalfont\small\bfseries}}
%% \makeatother

%% \makeatletter
%% \let\origsection\section
%% \renewcommand\section{\vspace{5em}\origsection}
%% \makeatother

% Don't print section numbers
\setcounter{secnumdepth}{0}

\setlength{\parindent}{0pt}
\setlength{\parskip}{1ex}

\pagestyle{fancy}
\newcommand{\makeheader}[3]{
  \fancyhead[L]{
    \begin{minipage}{30pt}
      \vspace{-3pt}
     \includegraphics[width=\textwidth]{logo.pdf}
    \end{minipage}
  }
  \fancyhead[C]{
   \LARGE{#1}
  }
  \fancyhead[R]{
   {\scriptsize
    #2 $\cdot$ Révision \#1 $\cdot$ \copyright\ \the\year$\quad$
    %% \begin{minipage}{30pt}
    %%  \vspace{-3pt}
    %%  \includegraphics[width=\textwidth]{logo.pdf} %defined at bottom
    %% \end{minipage}
   }}
  \fancyfoot[L]{}
  \fancyfoot[C]{}
  \fancyfoot[R]{}
}
\renewcommand\headrule{}

\newenvironment{absolutelynopagebreak}
  {\par\nobreak\vfil\penalty0\vfilneg
   \vtop\bgroup}
  {\par\xdef\tpd{\the\prevdepth}\egroup
    \prevdepth=\tpd}

\newcommand{\vb}{\usefont{T1}{DejaVuSansMono-TLF}{b}{n}}

\newenvironment{libcomments}{\begin{itemize}[topsep=-10000pt]}{\end{itemize}}

\setlength{\headheight}{20.1pt}
\setlength{\headsep}{10pt}

\renewcommand{\familydefault}{\sfdefault} % everything sans-serif

%Url handling
\hypersetup{
  colorlinks=true,
  urlcolor=black,
  linkcolor=black,
  citecolor=black
}

% font size -> set to equivalent to footnotesize of 10 pt
% it affect only the \\normalsize command, all are still relative to default's
% article one (10pt)
\newcommand\articlenormalsize{\fontsize{10pt}{12pt}\selectfont}
\renewcommand\normalsize{\footnotesize}

% redefine sections padding
\makeatletter
\renewcommand{\section}{\@startsection{section}{1}{0mm}%
                                {-1ex plus -.5ex minus -.2ex}%
                                {1ex}%
                                {\normalfont\articlenormalsize\bfseries}}
\renewcommand{\subsection}{\@startsection{subsection}{2}{4mm}%
                                {-1ex}%
                                {0.5ex}%
                                {\normalfont\articlenormalsize\bfseries}}
\renewcommand{\subsubsection}{\@startsection{subsubsection}{3}{16mm}%
                                {-1ex plus -.5ex minus -.2ex}%
                                {1ex plus .2ex}%
                                {\normalfont\small\bfseries}}
\makeatother

%\pagestyle{empty}
% set section counter to 0
\setcounter{secnumdepth}{0}

% Text at right
%\raggedright

% Set up the interline space in tables
\renewcommand{\arraystretch}{0.0}

% https://github.com/gpoore/minted/issues/82
\makeatletter
\newenvironment{catala}{%
  \VerbatimEnvironment
  \let\FV@ListVSpace\relax
  \begin{minted}[autogobble]{catala_fr}}%
 {\end{minted}}
\makeatother

\begin{document}

\makeheader{\MakeUppercase{La syntaxe de Catala} \hspace{2em}\small version française}{v0.8.0}{1}
% warning: this is repeated for p.2
% \begin{center}
%   2022-12-15
%   \hfill
%   {\Huge\bfseries\sffamily
%     Catala Syntax Cheat Sheet}
%   \hfill
%   v.0.8.0
% \end{center}
% \hrule
\null
\begin{multicols}{3}

\newlength\cola
\setlength\cola{\dimexpr .4\columnwidth -2\tabcolsep}
\newlength\colb
\setlength\colb{\dimexpr .6\columnwidth -2\tabcolsep}



\section{Programmation littéraire}

\begin{tabular}{@{}p{\cola}p{\colb}@{}}
  En-têtes & \begin{catala}
    # Titre
    ### Sous-sous-titre
  \end{catala}
  \\
  Référence au journal officiel & \begin{catala}
    # Article 1 | JORFARTI000012345678
    # Article 2 | LEGIARTI000012345678
    # Décision 3 | CETATEXT000012345678 
  \end{catala}
  \\
  Bloc de code &
  \begin{catala}
    ```catala
    ```
  \end{catala}
  \\
  Bloc de métadonnées & \begin{catala}
    ```catala-metadata
    ```
  \end{catala}
  \\
  Inclusion de fichier & \begin{catala}
    > Inclusion: foo.catala_en
  \end{catala}
  \\
\end{tabular}

\newcommand*\FancyVerbStartString{\PYG{l+s}{```catala}}
\newcommand*\FancyVerbStopString{\PYG{l+s}{```}}

\section{Types et littéraux}

\begin{tabular}{@{}p{\cola}p{\colb}@{}}
  \begin{catala}
    ```catala
    booléen
    ```
  \end{catala}
  &
  \begin{catala}
    ```catala
    vrai        faux
    ```
  \end{catala}
  \\
  \begin{catala}
    ```catala
    entier
    ```
  \end{catala}
  &
  \begin{catala}
    ```catala
    65536
    ```
  \end{catala}
  \\
  \begin{catala}
    ```catala
    décimal
    ```
  \end{catala}
  & \begin{catala}
      ```catala
      65536,262144        37%
      ```
    \end{catala}
  \\
  \begin{catala}
    ```catala
    argent
    ```
  \end{catala}
  &
  \begin{catala}
    ```catala
    1 234 567,89€
    ```
  \end{catala}
  \\
  \begin{catala}
    ```catala
    date
    ```
  \end{catala}
  &
  \begin{catala}
    ```catala
    |2021-01-31|
    ```
  \end{catala}
  \\
  \begin{catala}
    ```catala
    durée
    ```
  \end{catala}
  &
  \begin{catala}
    ```catala
    254 jour 4 mois 1 an
    ```
  \end{catala}
  \\
  \begin{catala}
    ```catala
    collection entier
    ```
  \end{catala}
  &
  \begin{catala}
    ```catala
    [ 12; 24; 36 ]
    ```
  \end{catala}
  \\
  \begin{catala}
    ```catala
    décimal dépend de argent
    ```
  \end{catala}
  &
  \begin{catala}
    ```catala
    f de x égale à x / 12,0€
    ```
  \end{catala}
  \\
  \begin{catala}
    ```catala
    Struct1
    ```
  \end{catala}
  &
  \begin{catala}
    ```catala
    Struct1 { -- chp1: 9 -- chp2: faux }
    ```
  \end{catala}
  \\
  \begin{catala}
    ```catala
    Énum1
    ```
  \end{catala}
  &
  \begin{catala}
    ```catala
    Cas1 contenu 12        Cas2
    ```
  \end{catala}
  \\
\end{tabular}

\section{Expressions}

\begin{tabular}{@{}p{\cola}p{\colb}@{}}
  Définition locale & \begin{catala}
    ```catala
    soit x égale à 36 - 5 dans ...
    ```
  \end{catala}
  \\
  Filtrage par motif & \begin{catala}
    ```catala
    selon expr sous forme
    -- Cas1 de x: ...
    -- Cas2 : ...
    ```
  \end{catala}
  \\
  Test de filtrage\newline avec variable optionnelle & \begin{catala}
    ```catala
    expr sous forme Cas1
    expr sous forme Cas1 de x
              et x >= 2
    ```
  \end{catala}
  \\
  Champ de structure & \begin{catala}
    ```catala
    struc1.chp2
    ```
  \end{catala}
  \\
  Appel de fonction & \begin{catala}
    ```catala
    f de 44,50€
    ```
  \end{catala}
  \\
  Var. de sous-ch. d'app. & \begin{catala}
    ```catala
    ss_ch1.var0
    ```
  \end{catala}
  \\
  Appel direct\newline du champ d'application & \begin{catala}
    ```catala
    résultat de ChApp1
      avec { -- chp1: 9 -- chp2: vrai }
    ```
  \end{catala}
  \\
  Branchement conditionnel & \begin{catala}
    ```catala
    si ... alors ... sinon ...
    ```
  \end{catala}
  \\
\end{tabular}

\section{Déclaration des métadonnées}

\begin{tabular}{@{}p{\cola}p{\colb}@{}}

  Déclaration de structure &
  \begin{catala}
    ```catala
    déclaration structure Struct1:
      donnée chp1 contenu entier
      donnée chp2 contenu booléen
  \end{catala}
  \\
  Déclaration d'énumération & \begin{catala}
    ```catala
    déclaration énumeration Énum1:
      -- Cas1 contenu entier
      -- Cas2
    ```
  \end{catala}
  \\
  Déclaration de champ\newline d'application & \begin{catala}
    ```catala
    déclaration champ d'application ChApp1:
      interne var1 contenu entier
      interne var2 condition
      ss_ch1 champ d'application ChApp0
    ```
  \end{catala}
  \\
  Qualificateurs d'entrée-\newline sortie & \begin{catala}
    ```catala
      interne var1 contenu ...
      résultat var3 contenu ...
      entrée var4 contenu ...
      entrée résultat var5 contenu ...
      contexte var6 contenu ...
      contexte résultat var7 contenu ...
    ```
  \end{catala}
  \\
  Transitions d'état & \begin{catala}
    ```catala
    interne var1 contenu ...
      état avant
      état après
    ```
  \end{catala}
  \\
\end{tabular}

\section{Opérations}

\begin{tabular}{@{}p{\cola}p{\colb}@{}}
  Opérateurs logiques & \begin{catala}
    ```catala
    non a
    a et b
    a ou b        # "ou à défaut"
    a ou bien b   # ou exclusif
    ```
  \end{catala}
  \\
  Arithmétique & \begin{catala}
    ```catala
    - a
    a + b        a - b
    a * b        a / b
    ```
  \end{catala}
  \\
  Comparaisons &
  \begin{catala}
    ```catala
    a = b        a != b
    a > b        a < b
    a >= b       a <= b
    ```
  \end{catala}
  \\
  Conversions & \begin{catala}
    ```catala
    décimal de 44
    argent de 23,15
    ```
  \end{catala}
  \\
  Arrondis & \begin{catala}
    ```catala
    arrondi de 9,99€
    ```
  \end{catala}
  \\
  Éléments de dates & \begin{catala}
    ```catala
    accès_jour de ...
    accès_mois de ...
    accès_année de ...
    ```
  \end{catala}
  \\
  Opérateurs à types\newline explicites &
  \begin{catala}
    ```catala
    a +! b    # entier
    a +. b    # décimal
    a +€ b    # argent
    a +^ b    # durée
    ```
  \end{catala}
  \\
\end{tabular}

\columnbreak

\section{Définition de champ d'application}

\begin{tabular}{@{}p{\cola}p{\colb}@{}}
  Utilisation & \begin{catala}
    ```catala
    champ d'application ChApp1: ...
    ```
  \end{catala}
  \\
  Avec condition générale & \begin{catala}
    ```catala
    champ d'application ChApp1
      sous condition var1 >= 2: ...
    ```
  \end{catala}
  \\
  Définition inconditionnelle & \begin{catala}
    ```catala
    définition var1 égale à ...
    ```
  \end{catala}
  \\
  Définition conditionnelle& \begin{catala}
    ```catala
    définition var1
      sous condition ...
      conséquence égale à ...
    ```
  \end{catala}
  \\
  Règle\newline(définition de condition) & \begin{catala}
    ```catala
    règle var2
      sous condition var1 >= 2
      conséquence rempli
    ```
  \end{catala}
  \\
  Règle négative & \begin{catala}
    ```catala
      conséquence non rempli
    ```
  \end{catala}
  \\
  Définition de fonction/règle & \begin{catala}
    ```catala
    définition f de x égale à ...
    ```
  \end{catala}
  \\
  Définition/règle étiquetée & \begin{catala}
    ```catala
    étiquette étq1 définition var1 ...
    ```
  \end{catala}
  \\
  Exc. à définition étiquetée & \begin{catala}
    ```catala
    exception étq1 définition var1 ...
    ```
  \end{catala}
  \\
  Exception à implicite & \begin{catala}
    ```catala
    exception définition var1 ...
    ```
  \end{catala}
  \\
  Définition d'états & \begin{catala}
    ```catala
    définition var1
      état avant
      égale à ...
    ```
  \end{catala}
  \\
  Assertion & \begin{catala}
    ```catala
    assertion ...
    ```
  \end{catala}
  \\
\end{tabular}

\section{Opérations sur les collections}

\begin{tabular}{@{}p{\cola}p{\colb}@{}}
  
  Test de présence& \begin{catala}
    ```catala
    coll contient 3
    ```
  \end{catala}
  \\
  Cardinal & \begin{catala}
    ```catala
    nombre de coll
    ```
  \end{catala}
  \\
  Test d'existence & \begin{catala}
    ```catala
    existe x parmi coll tel que x >= 2
    ```
  \end{catala}
  \\
  Test pour tous & \begin{catala}
    ```catala
    pour tout x parmi coll on a x >= 2
    ```
  \end{catala}
  \\
  Application un-à-un & \begin{catala}
    ```catala
    (x + 2) pour x parmi coll
    ```
  \end{catala}
  \\
  Filtrage & \begin{catala}
    ```catala
    x parmi coll tel que x >= 2
    ```
  \end{catala}
  \\
  Filtrage + application & \begin{catala}
    ```catala
    (x - 2) pour x parmi coll
            tel que x >= 2
    ```
  \end{catala}
  \\
  Réunion & \begin{catala}
    ```catala
    coll1 ++ coll2
    ```
  \end{catala}
  \\
  Aggrégation & \begin{catala}
    ```catala
    somme entier coll
    ```
  \end{catala}
  \\
  Comptage & \begin{catala}
    ```catala
    nombre de coll
    ```
  \end{catala}
  \\
  Extremums & \begin{catala}
    ```catala
    maximum de coll
      ou si collection vide alors -1
    ```
  \end{catala}
  \\
  Élément selon extremum & \begin{catala}
    ```catala
    x parmi coll
      tel que (x * x) est minimum
      ou si collection vide alors -1
    ```
  \end{catala}
  \\
\end{tabular}

\end{multicols}
\end{document}
